% Generated draft: requires mathematical and source review.
\documentclass[11pt]{article}
\usepackage[margin=25mm]{geometry}
\usepackage{fontspec}
\setmainfont{Latin Modern Roman}
\usepackage{amsmath,amssymb,mathtools,unicode-math}
\setmathfont{Latin Modern Math}
\usepackage{enumitem,needspace,xurl,hyperref}
\setlength{\parindent}{0pt}
\setlength{\parskip}{5pt}
\setlength{\emergencystretch}{3em}
\newcommand{\N}{\mathbb{N}}
\newcommand{\Z}{\mathbb{Z}}
\newcommand{\Q}{\mathbb{Q}}
\newcommand{\R}{\mathbb{R}}
\newcommand{\C}{\mathbb{C}}
\newcommand{\F}{\mathbb{F}}
\newcommand{\A}{\mathbb{A}}
\newcommand{\PP}{\mathbb{P}}
\newcommand{\E}{\mathbb{E}}
\newcommand{\eps}{\varepsilon}
\newcommand{\dd}{\,\mathrm{d}}
\begin{document}
\section*{UnsolvedMath ID: 20000242}
\textbf{Source problem code:} catalogue entry
\subsection{Definitions and mathematical statement}
For a finitely generated module $M$ over a ring $R$, let $\mu_R(M)$ be its minimum number of generators. The main question, attributed in the source to Satya Mandal, is whether an ideal $I$ in $R=k[X_1,\ldots,X_n]$ satisfies
\[
 \mu_R(I)=\mu_{R/I}(I/I^2).
\]
The field is unspecified. The usual efficient-generation problem concerns proper ideals; the extracted statement does not say proper, and $I=R$ would give different sides. That missing qualification is recorded rather than silently removed. The source reports the equality known when $\mu(I/I^2)\geq\dim(R/I)+2$.

The related Bhatwadekar question takes a surjective $\C$-algebra map $\alpha:\C[X_1,\ldots,X_5]\twoheadrightarrow\C[Y_1,Y_2]$ and $I=\ker\alpha$. It reports $\mu(I/I^2)=3$ and asks whether $\mu(I)=3$. Asok's remark suggests detecting the number of generators of a projective module through lifting its classifying map $\operatorname{Spec}R\to BGL_r$ to a suitable Grassmannian, without giving the required hypotheses. Fasel asks for a Chow-type computable cohomological obstruction to a smooth embedding $Y\subset X=\operatorname{Spec}R$ being a global complete intersection, and subsequently for a locally complete-intersection $Y$.

Fasel's universal setup in the damaged remarks is
\[
 R_{\mathrm{univ}}=k[X_1,\ldots,X_n,Y_1,\ldots,Y_n,Z]/
       (\textstyle\sum_{i=1}^nX_iY_i-Z(1+Z)),
\]
\[
 X_{\mathrm{univ}}=\operatorname{Spec}R_{\mathrm{univ}},\qquad
 Y_{\mathrm{univ}}=\operatorname{Spec}(R_{\mathrm{univ}}/(X_1,\ldots,X_n,Z)).
\]
It describes $X_{\mathrm{univ}}$ as an $\A^1$-sphere and asserts a classifying map of pairs $(X,Y)\to(X_{\mathrm{univ}},Y_{\mathrm{univ}})$ from conormal-generation data, but the value of the conormal-generator bound is lost. It reports a Chow--Witt decomposition $GW(k)\oplus GW(k)\alpha$, with $\alpha$ a form on $Y_{\mathrm{univ}}$, while its precise grading and source twist are damaged. The related question asks whether $f^*(\alpha)=0$ in the indicated Chow--Witt group of $X$ twisted by $\omega_{X/k}$ suffices for $Y$ to be a global complete intersection. Choices of orientation, grading, and classifying map must be recovered before this is a precise obstruction criterion.

The research aid reports the general and complex-kernel questions open and proposes a restricted strong-lifting criterion. For Segre data
\[
 I=(a_1,\ldots,a_r,s),\qquad s(1-s)=\sum_{i=1}^ra_ib_i,
\]
if the opposite-coordinate column $b=(b_i)$ can be completed to the first basis vector by an invertible change of basis, it claims the prescribed conormal orientation lifts to $r$ generators of $I$. A strong lift here means generators congruent to the prescribed $a_i$ modulo $I^2$. Over a polynomial ring over a field, it reports that $(b_1,\ldots,b_r)=R$ suffices by Quillen--Suslin. Its contrapositive candidate is that every nonliftable orientation has a proper opposite-coordinate ideal in every Segre presentation. This conditional locus is not a solution of the unrestricted efficient-generation or cohomological questions.
\subsection{Short English statement}
Can an ideal in a polynomial ring be generated by as few elements as its conormal module, including the stated complex kernel example? The source also asks for homotopical/cohomological complete-intersection obstructions. The report gives only a special strong-lifting criterion, while substantial OCR damage leaves some universal obstruction data unresolved.
\subsection{Sources}
\begin{itemize}
\item Local UnsolvedMath registry, permanent ID $20000242$, source code \texttt{catalogue entry}; created and updated 2026-08-14. Status \texttt{partially\_solved}; classification \texttt{PARTIAL-PROGRESS}; difficulty L3; proposer field null.
\item AIM Problem Lists, Algebraic geometry, workshop \emph{Projective modules and $\A^1$-homotopy theory}, May 5--9, 2014, Question 5, source type PDF, \url{https://aimath.org/pastworkshops/projectiveA1problems.pdf}. Canonical location \texttt{aim-algebraic-geometry-notes.json}, \texttt{notes[241]}, tag \texttt{question}; statement status \texttt{exact}. Verification reports checking substantial OCR damage against the workshop PDF; this draft makes no new source-visit claim. Separate literature field empty.
\item Main attribution Satya Mandal; related questions/remarks S. M. Bhatwadekar, Aravind Asok, and Jean Fasel. The research summary calls the main problem Murthy's efficient-generation question and invokes Quillen--Suslin. No exact primary bibliography is supplied for these results or the universal Chow--Witt assertion.
\item Supplied research attempt 1: result \texttt{partial\_result}, run status \texttt{open}, rigor label \texttt{proved}, low novelty confidence. Candidate contribution \texttt{proposition}: completable opposite-coordinate strong lifting and its proper-ideal contrapositive. Explicit caveat: machine-generated research aid, not peer reviewed; independently verify before citation.
\item Dataset attribution: Ulam AI / UnsolvedMath Contributors, CC BY 4.0; original problem attribution remains to the supplied AIM workshop collection.
\end{itemize}
\end{document}
